\begin{abstract}
The strongest existing LLM-agent security benchmarks---Agent Security Bench, AgentDojo,
InjecAgent, ToolEmu---are single-agent, yet deployed systems are collectives in which a
planner directs workers, agents share memory, and tools are third-party. The phenomena
that dominate risk there---collusion, a compromised orchestrator, malicious tools, Sybils,
group-level emergent failure---cannot be produced by a single-agent harness. We present
\defname{}, an environment and evaluation methodology for multi-agent LLM security: a
forced-coordination task suite with heterogeneous roles and partial observability over star,
chain, tree, and mesh topologies at $3$--$50$ agents; a parameterized adversary model; a
system-level metric $\mathrm{NRP}_{\mathrm{sys}}=\mathrm{PNA}_{\mathrm{sys}}(1-\mathrm{ASR}_{\mathrm{sys}})$
scored by a deterministic checker over the global state trace, never an LLM judge on the
critical path; and \corb, a co-evolving red/blue protocol logging attack--defense
trajectories and a (security, cost) Pareto frontier. Supporting theory, proved under
explicit assumptions, characterizes the metric axiomatically (reducing to ASB at one agent),
bounds compromise propagation across topologies, establishes a seed-placement advantage and
a forced-coordination lower bound, and proves evaluator integrity. On a clearly-labelled
\emph{synthetic} emulator, a per-agent guardrail's relative reduction in action-channel
attack success rate shrinks $41.6\%\to31.8\%\to27.6\%$ from single-agent to collusion to
compromised-orchestrator context. We then ask what that guardrail buys on four real
backbones in three model families: a prompt-level defense drives action-channel compliance
to $0/1075$ decisions, so action-channel attack success is $0$ throughout, yet
system-level attack success remains $1.000$ on tree and star, because once the unsafe
action is blocked the system predicate is still discharged by deterministic availability
sabotage that topology alone determines. A fourth backbone ignores the defense entirely. A
defense measured on the action channel therefore cannot certify system-level security, and
its efficacy is a property of the model as much as of the defense.
\end{abstract}